\section{Diseño e implementación en Noir}
\label{sec:noir_impl}

La implementación en Noir expresa los dos enunciados presentados en la
Sección~\ref{sec:puente_arquitectura} mediante tres circuitos, dos para el enunciado de consenso y
uno para el de inclusión, cuyas pruebas genera Barretenberg con el sistema UltraHonk. Comparte con
la implementación en SP1 el token, el contrato \texttt{Sender}, la interfaz y el recorrido general
de una transferencia, pero difiere en la organización de las pruebas y en los contratos de
destino, diferencias que resume la Sección~\ref{sec:puente_variantes}. Esta sección describe lo
que establece cada circuito, las consecuencias de expresarlo como un conjunto de restricciones,
la codificación de sus valores públicos, los contratos que los consumen y la integración con el
\textit{relayer}. La descripción corresponde al código de los repositorios del puente y de sus
circuitos~\cite{guaranibridgenoir, zkbridgezerorepo, noirmerklerepo}, con las versiones de las
herramientas indicadas en la Sección~\ref{sec:noir}, Nargo~1.0.0-beta.22 y
Barretenberg~5.0.0-nightly.20260522.

\subsection{Programa de consenso}
\label{sec:noir_impl_consenso}

El enunciado de consenso se reparte entre dos circuitos. Expresado en un único circuito, con
unos 13,3 millones de compuertas, superó los 91 GB de memoria durante la compilación en el equipo
disponible, y el sistema operativo terminó el proceso (Sección~\ref{sec:resultados_etapas}). El
primero, denominado circuito del comité, recibe como entradas privadas las 512 claves públicas
comprimidas del comité, el vector de participación, la firma agregada, los cinco campos del
encabezado de la cadena de baliza firmado, el dominio de firma, el hash del bloque de ejecución
junto con la rama SSZ que lo vincula con \texttt{body\_root} y el encabezado de ejecución
codificado en RLP (Sección~\ref{sec:eth_consenso}). El circuito comprueba que las claves
corresponden al compromiso del comité, que el número de participantes multiplicado por tres
alcanza al menos 1024, es decir, que firmaron 342 o más de los 512 miembros, y suma las claves de
los participantes. Deriva después la raíz de firma como el hash SHA-256 de la raíz del encabezado,
calculada mediante \texttt{hash\_tree\_root}~\cite{sszspec}, concatenada con el
dominio~\cite{ethconsensusphase0}, y la transforma en un punto de $\mathbb{G}_2$ según el
procedimiento de la RFC~9380, con la etiqueta de separación de dominio que emplea
Ethereum~\cite{rfc9380}. Comprueba, además, que la firma pertenece a la curva y al subgrupo
$\mathbb{G}_2$, que la rama SSZ conduce desde el hash del bloque de ejecución hasta
\texttt{body\_root} y que el hash Keccak-256 del encabezado de ejecución coincide con ese hash. De
ese encabezado extrae, por último, las raíces de transacciones y de recibos
(Tabla~\ref{tab:eth_header}), que expone como valores públicos.

El segundo, denominado circuito de emparejamiento, recibe la clave agregada, el punto del mensaje
y la firma, y verifica la Ecuación~\eqref{eq:bls_agregada}. Los dos emparejamientos de esa
ecuación comparten el bucle de Miller y una única exponenciación final, lo que reduce el tamaño
del circuito respecto de evaluarlos por separado. Ninguno de los dos circuitos expone los tres
puntos, que ocupan varios cientos de bytes, sino un valor de enlace calculado por ambos como el
hash Poseidon2 de la clave agregada, el punto del mensaje y la firma. \texttt{BeaconLightClient}
exige que ese valor coincida en las dos pruebas, de modo que una prueba de emparejamiento
obtenida para otra firma no puede combinarse con la prueba del comité.

Dos decisiones de diseño determinan la utilidad de estas pruebas para el puente. La primera
consiste en derivar el mensaje firmado dentro del circuito, en lugar de recibirlo como un valor
ya calculado. Un mensaje recibido de ese modo sería opaco, ya que la prueba no permitiría saber de
qué encabezado proviene ni qué bloque de ejecución representa, y esa era precisamente la
limitación de la versión anterior del circuito, que recibía el punto del mensaje y la clave
pública como entradas libres (Sección~\ref{sec:metodologia_alcance}). La segunda consiste en sumar
las claves dentro del circuito y exponer el compromiso del comité completo en lugar de la clave
agregada. La clave agregada depende de qué miembros firmaron y cambia de un \textit{slot} a otro,
mientras que el compromiso permanece estable durante el período del comité, de 8192
\textit{slots} o unas 27 horas (Sección~\ref{sec:eth_sync_committee}), lo que permite que el
propietario de \texttt{BeaconLightClient} lo registre una vez por período. El compromiso se calcula
con Poseidon2, una función de hash algebraica cuyo costo como circuito es muy inferior al de
SHA-256 sobre los 24\,576 bytes de las claves, a diferencia de SP1, que emplea SHA-256
(Sección~\ref{sec:sp1_impl_consenso}).

La expresión de este enunciado en Noir se ve favorecida porque casi todos sus datos tienen un
tamaño fijado por el protocolo. El comité tiene siempre 512 miembros, cada clave comprimida ocupa
48 bytes y la firma 96, y la rama SSZ tiene once niveles, correspondientes al índice generalizado
2856 con que el cuerpo del bloque de baliza localiza el hash del bloque de ejecución en la
bifurcación vigente en Sepolia durante las pruebas. Los arreglos de tamaño fijo que exige el
programa principal (Sección~\ref{sec:noir_descripcion}) no imponen, por ello, máximos arbitrarios,
aunque una bifurcación que modifique el dominio, la posición de ese campo o el formato del
encabezado de ejecución obliga a regenerar el circuito del comité y su verificador. El costo se
concentra en la aritmética de la curva BLS12-381. Noir no ofrece funciones de caja negra para esa
curva~\cite{noirdocsbeta}, de modo que la descompresión de las claves, que exige calcular raíces
cuadradas en el cuerpo base de 381 bits, su suma, la transformación del mensaje en un punto de
$\mathbb{G}_2$ y los emparejamientos se expresan mediante aritmética no nativa, con varios
elementos del cuerpo nativo para cada elemento de ese cuerpo (Sección~\ref{sec:noir_limitaciones}).
En un circuito, además, una instrucción condicional no evita el costo de la rama que no se
ejecuta, ya que el compilador genera restricciones para ambas~\cite{noirdocsbeta}, y la suma
condicionada por el vector de participación cuesta lo mismo que la del comité completo, con
independencia del número de firmantes. Los hashes SHA-256 del enunciado, entre ellos los de la
raíz del encabezado, la rama y la expansión del mensaje, se apoyan en la función de compresión que
Noir ofrece como caja negra. Cuando una operación costosa, como una raíz cuadrada, se calcula
mediante una función no restringida, su resultado solo queda respaldado por la prueba si el
programa lo restringe explícitamente, conforme al patrón de la Figura~\ref{fig:noir_unconstrained}.

\subsection{Programa de inclusión}
\label{sec:noir_impl_inclusion}

El circuito de inclusión recibe las raíces de transacciones y de recibos, la posición de la
transacción de bloqueo, la transacción y su recibo serializados y los nodos de las ramas del
árbol de Merkle-Patricia que conducen desde cada raíz hasta la hoja de esa posición
(Sección~\ref{sec:eth_mpt}). A diferencia del programa de SP1, no recibe el encabezado de
ejecución ni las demás transacciones del bloque, ya que la vinculación de las raíces con el
encabezado la establece el circuito del comité, y no reconstruye los árboles completos, sino que
verifica una rama en cada uno. Cada rama tiene hasta seis nodos, uno por nivel. El circuito
comprueba que el hash Keccak-256 del primer nodo coincide con la raíz y que el de cada nodo
siguiente coincide con la referencia que guarda su padre, recorre los nodos de ramificación y de
extensión según los \textit{nibbles} de la clave, que es la codificación RLP de la posición, y
exige que la hoja contenga exactamente la transacción o el recibo serializado. El procedimiento se
aplica dos veces con la misma posición, una en el árbol de transacciones y otra en el de recibos,
de modo que no es posible presentar el recibo exitoso de otra transacción. Como todas las
referencias entre nodos deben ser hashes, el circuito no admite los nodos de menos de 32 bytes
que el árbol puede incrustar en su padre.
Decodifica después el sobre de la transacción, cuyo formato depende de su
tipo~\cite{eip2718}, y exige que se dirija directamente al contrato \texttt{Sender} y que sus datos
correspondan a una llamada a \lstinline|lock(address,uint256)|, identificada por el selector
\texttt{0x282d3fdf}, de la que extrae el destinatario y el monto. Se admiten las transacciones
heredadas y las de los tipos 1 a 4. Del recibo lee el código de estado, que debe indicar una
ejecución exitosa~\cite{eip658}, y busca entre sus eventos uno \texttt{Locked} emitido por
\texttt{Sender}, del que obtiene el identificador del bloqueo, el destinatario y el monto. Los
dos últimos deben coincidir con los de la llamada.

La comprobación del evento no es necesaria para la seguridad del puente, puesto que una llamada
directa a \texttt{lock} que termina sin revertir implica que el evento se emitió
(Sección~\ref{sec:puente_aceptacion}). Cumple, sin embargo, dos funciones en esta implementación.
Aporta el identificador del bloqueo, que el contrato asigna durante la ejecución y que no figura en
los datos de la llamada, y que \texttt{ReceiverZKBound} emplea para impedir que una operación se
procese dos veces, y añade un control redundante frente a errores en la decodificación de la
llamada. El programa de inclusión de SP1, en cambio, no examina los eventos del recibo y obtiene
el destinatario y el monto únicamente de los datos de la llamada
(Sección~\ref{sec:sp1_impl_inclusion}). En ambos casos, la exigencia de una llamada directa
implica que un bloqueo realizado a través de otro contrato no puede demostrarse. Esta limitación
tiene consecuencias prácticas, porque las billeteras comienzan a convertir las cuentas de sus
usuarios en cuentas inteligentes mediante la delegación de código que introduce la
EIP-7702~\cite{eip7702}, con la que el bloqueo llega a \texttt{Sender} como una llamada interna.
Las pruebas en Optimism Sepolia registraron un caso de este tipo (Sección~\ref{sec:resultados_gas}).

La expresión de este enunciado como circuito presenta una dificultad que no aparece en el de
consenso, porque la transacción, el recibo y los nodos de las ramas no tienen un tamaño fijado por
el protocolo. El recibo, en particular, incluye todos los eventos emitidos. Como los parámetros
del programa principal deben tener un tamaño fijo, el circuito recibe arreglos dimensionados para
una longitud máxima, acompañados de las longitudes efectivas
(Sección~\ref{sec:noir_limitaciones}). Las operaciones cuyo recorrido depende de los datos, como
la decodificación RLP o el cálculo de Keccak-256 sobre entradas de longitud variable a partir de
la permutación que Noir ofrece como caja negra, se expresan mediante ciclos acotados por esos
máximos, que el compilador desenrolla, y mediante accesos a posiciones de los arreglos que solo se
conocen durante la ejecución, más costosos que los accesos a posiciones fijas~\cite{noirdocsbeta}.
El circuito incorpora, por ello, las restricciones correspondientes a los tamaños máximos con
independencia de la transacción concreta. Para el árbol de transacciones se adoptaron cotas
menores que para el de recibos, con hojas de hasta 256 bytes y nodos de hasta 543, lo que basta
para una transacción de bloqueo, que ocupa entre 173 y 181 bytes. En el árbol de recibos, las
hojas admiten hasta 784 bytes y los nodos hasta 800. Con esas cotas el circuito tiene
4\,105\,601 compuertas y cabe en $2^{22}$ filas, mientras que con las cotas de los recibos
aplicadas a ambos árboles alcanzaba unos 4,94 millones de compuertas y requería $2^{23}$.
Una transacción que supere esos máximos no puede demostrarse, y aumentarlos modifica el circuito
y, con él, su clave de verificación y su contrato verificador.

\subsection{Valores públicos y contratos verificadores}
\label{sec:noir_impl_publicas}

En Noir, los valores públicos se expresan como entradas públicas del circuito, cada una de las
cuales es un elemento del cuerpo de 254 bits de Barretenberg (Sección~\ref{sec:noir_descripcion}).
Un valor de 32 bytes, como un hash, puede superar el módulo de ese cuerpo y debe descomponerse en
varios elementos, ya sea en sus bytes, como ocurre con los arreglos de tipo \lstinline|u8|
(Sección~\ref{sec:noir_limitaciones}), o en dos mitades de 128 bits. El circuito del comité adopta
la segunda opción para las dos raíces y expone siete entradas, el circuito de emparejamiento
expone solo el valor de enlace y el de inclusión expone 168 entradas, una por byte. La
Tabla~\ref{tab:valores_publicos_noir} detalla su contenido y el uso que hacen de ellas los
contratos de destino. La descomposición y el orden de las entradas forman parte de la interfaz
entre el circuito y el contrato, y una discrepancia entre ambos haría que el contrato comparara
valores distintos de los demostrados. La descomposición incide, además, en el costo, ya que el
verificador generado incorpora cada entrada pública en la transcripción de Fiat-Shamir y en el
término que la vincula con la prueba, de modo que el gas crece con el número de
entradas~\cite{barretenbergsrc}, efecto que se cuantifica en la Sección~\ref{sec:resultados_gas}.

\begin{table}[h!]
	\centering
	\caption{Entradas públicas de los circuitos de Noir y uso que hacen de ellas los contratos de
	destino.}
	\label{tab:valores_publicos_noir}
	\footnotesize
	\begin{tabular}{>{\raggedright\arraybackslash}p{3.7cm} >{\raggedright\arraybackslash}p{5.6cm}
		>{\raggedright\arraybackslash}p{6.2cm}}
		\toprule
		\textbf{Entrada} & \textbf{Contenido} & \textbf{Uso en el contrato} \\
		\midrule
		\multicolumn{3}{l}{\textit{Circuito del comité, siete elementos, consumido por}
		\texttt{BeaconLightClient}} \\
		\addlinespace
		\texttt{committee\_commitment} & Hash Poseidon2 de las 512 claves del comité. & Debe
		coincidir con el registrado por el propietario para el período del \textit{slot}. \\
		\texttt{slot} & \textit{Slot} del encabezado firmado. & Determina el período, igual al
		cociente entre el \textit{slot} y 8192. \\
		\texttt{link} & Hash Poseidon2 de la clave agregada, el punto del mensaje y la firma. & Debe
		coincidir con el del circuito de emparejamiento. \\
		\texttt{transactions\_root} (dos mitades) & Raíz de transacciones del bloque de ejecución. &
		Se registra junto con la raíz de recibos. \\
		\texttt{receipts\_root} (dos mitades) & Raíz de recibos del mismo bloque. & Se registra junto
		con la raíz de transacciones. \\
		\midrule
		\multicolumn{3}{l}{\textit{Circuito de emparejamiento, un elemento, consumido por}
		\texttt{BeaconLightClient}} \\
		\addlinespace
		\texttt{link} & Mismo valor de enlace. & Debe coincidir con el del circuito del comité. \\
		\midrule
		\multicolumn{3}{l}{\textit{Circuito de inclusión, 168 elementos de un byte, consumido por}
		\texttt{ReceiverZKBound}} \\
		\addlinespace
		\texttt{receiptsRoot} (32) & Raíz de recibos. & El par de raíces debe estar registrado. \\
		\texttt{sender} (20) & Contrato que emitió el evento y recibió la llamada. & Debe ser el
		\texttt{Sender} configurado. \\
		\texttt{id} (32) & Identificador del bloqueo, tomado del evento. & Debe coincidir con el
		argumento de la llamada y no haberse procesado antes. \\
		\texttt{recipient} (20) & Destinatario en la cadena de destino. & Debe coincidir con el
		argumento de la llamada. \\
		\texttt{amount} (32) & Monto bloqueado. & Debe coincidir con el argumento de la llamada. \\
		\texttt{transactionsRoot} (32) & Raíz de transacciones. & El par de raíces debe estar
		registrado. \\
		\bottomrule
	\end{tabular}
\end{table}

\texttt{BeaconLightClient} recibe las pruebas de consenso mediante su función
\lstinline|update|. Comprueba que el circuito del comité aporta siete entradas y el de
emparejamiento una, que el compromiso del comité coincide con el registrado para el período del
\textit{slot} y que el valor de enlace es el mismo en ambas pruebas, delega la verificación de cada
una en su contrato verificador y, si todas las comprobaciones se cumplen, registra como conocido
el hash Keccak-256 de la concatenación de las dos raíces y emite un evento con el \textit{slot} y
ambas raíces. \texttt{ReceiverZKBound} acuña mediante su función \lstinline|release|, que recibe
la prueba de inclusión, sus entradas públicas y el identificador, el destinatario y el monto de la
operación. Exige que el par de raíces esté registrado en \texttt{BeaconLightClient}, que las
entradas coincidan con esas raíces, con el \texttt{Sender} configurado y con los argumentos de la
llamada, que la prueba sea válida y que el identificador no se haya procesado antes. Registrar el
par, y no cada raíz por separado, impide combinar la transacción de un bloque con el recibo de
otro.

El contrato verificador de cada circuito lo genera Barretenberg a partir de su clave de
verificación, con el destino \lstinline|evm| descrito en la Sección~\ref{sec:noir_flujo}. El
contrato generado incorpora como constantes la clave de verificación, que incluye el tamaño del
circuito y el número de entradas públicas, y expone una función \lstinline|verify| que recibe la
prueba y el arreglo de entradas. Esta función comprueba la longitud de la prueba y el número de
entradas, recalcula los desafíos de Fiat-Shamir con Keccak, verifica el protocolo
\textit{sumcheck} y la apertura de los compromisos KZG mediante los contratos precompilados de la
curva BN254, y revierte la transacción si alguna comprobación falla~\cite{barretenbergsrc}. Como
la clave forma parte del contrato, cada circuito requiere su propio verificador, de modo que el
puente despliega tres, \texttt{CommitteeVerifier}, \texttt{PairingVerifier} y
\texttt{TxInclusionVerifier}, cuyo código ocupa entre 17\,155 y 17\,532 bytes, por debajo del
límite de 24\,576 bytes que establece la EIP-170~\cite{eip170}. Fijar en los contratos de destino
la dirección de cada verificador cumple la misma función que fijar la clave del programa en SP1,
ya que un verificador distinto aceptaría pruebas de otro circuito. Los tres emplean la cadena de
referencia universal de la ceremonia Ignition, por lo que no requieren una ceremonia específica
para cada circuito (Sección~\ref{sec:noir_flujo}).

\subsection{Generación de las pruebas e integración}
\label{sec:noir_impl_integracion}

La preparación se realiza una vez por circuito. Nargo compila los tres programas, Barretenberg
obtiene la clave de verificación de cada circuito y genera su contrato verificador, y los
verificadores se despliegan en la cadena de destino antes que \texttt{BeaconLightClient} y
\texttt{ReceiverZKBound}, en los que se configuran sus direcciones. La compilación del circuito de
emparejamiento, el mayor de los tres, es la etapa más exigente de esta preparación
(Sección~\ref{sec:resultados_etapas}). Cualquier modificación de un programa, incluida la de sus
tamaños máximos, exige repetir estos pasos para el circuito afectado. Al inicio de cada período,
el propietario registra el compromiso del nuevo comité mediante la función
\lstinline|setCommittee|.

Durante cada transferencia, el \textit{relayer} espera el encabezado de la cadena de baliza
construido sobre el bloque del bloqueo y obtiene del nodo de baliza la firma del comité sobre ese
encabezado, que el nodo ofrece durante un intervalo del orden de un \textit{slot}. Si ese intervalo
transcurre sin que la firma se obtenga, el bloque no puede anclarse con este procedimiento. A
continuación organiza los datos según la interfaz de cada circuito, lo que incluye completar los
arreglos hasta sus tamaños máximos e indicar las longitudes efectivas, calcula los testigos y
genera las tres pruebas con el ejecutable nativo \texttt{bb} y el destino \lstinline|evm|. La
versión de Barretenberg para WebAssembly, \texttt{bb.js}, no pudo generar la prueba de inclusión,
porque excede el límite de memoria de 32 bits de ese entorno. Si el par de raíces del bloque no
está registrado, el \textit{relayer} envía las dos pruebas de consenso a \lstinline|update| en una
transacción, y envía después la de inclusión a \lstinline|release| en otra. El anclaje se paga,
por lo tanto, una vez por bloque y lo aprovechan todas las transferencias de ese bloque, mientras
que la acuñación se paga en cada transferencia. Una ejecución que no satisface las aserciones de
un circuito no produce testigo, de modo que esos fallos se detectan antes de generar la prueba.

Las etapas de preparación, de puesta en marcha y de cada operación distinguidas en la
Sección~\ref{sec:noir_flujo} se aplican del mismo modo a los tres circuitos, y su medición se
define en la metodología de comparación. La sección siguiente describe la expresión de los mismos
enunciados en SP1.
